\lysh{33696_SOLUTION}{1}
% TODO: TikZ σχήμα
(Α) Έστω $K(x_K, y_K)$ το κέντρο του κύκλου.
Αφού το $K$ είναι σημείο της ευθείας $\varepsilon: y = 2x - 1$, ισχύει $y_K = 2x_K - 1$. Άρα $K(x_K, 2x_K - 1)$.
Ο κύκλος $\mathcal{C}$ εφάπτεται της ευθείας $\zeta: x + y - 2 = 0$ άρα ισχύει $d(K, \zeta) = \rho$.
Έχουμε:
\[
d(K, \zeta) = \rho \iff \frac{|x_K + (2x_K - 1) - 2|}{\sqrt{1^2 + 1^2}} = 3\sqrt{2} \iff \frac{|3x_K - 3|}{\sqrt{2}} = 3\sqrt{2}
\]
\[
\iff |3x_K - 3| = 3\sqrt{2} \cdot \sqrt{2} \iff |3x_K - 3| = 6
\]
\[
\iff 3x_K - 3 = 6 \quad \text{ή} \quad 3x_K - 3 = -6 \iff 3x_K = 9 \quad \text{ή} \quad 3x_K = -3
\]
\[
\iff x_K = 3 \quad \text{ή} \quad x_K = -1.
\]
Αφού το $K$ είναι σημείο του πρώτου τεταρτημορίου είναι $x_K > 0$ οπότε $x_K = 3$.
Τότε $y_K = 2(3) - 1 = 5$. Επομένως το κέντρο του κύκλου είναι το $K(3,5)$ και η εξίσωση του κύκλου είναι:
\[
\mathcal{C}: (x - 3)^2 + (y - 5)^2 = (3\sqrt{2})^2
\]
\[
\mathcal{C}: (x - 3)^2 + (y - 5)^2 = 18.
\]

(β)
i. Η ευθεία $KA$ είναι κάθετη στην εφαπτομένη $\zeta$ του κύκλου $\mathcal{C}$ στο $A$, άρα ισχύει:
\[
\lambda_{KA}\lambda_{\zeta} = -1.
\]
Η ευθεία $\zeta: x + y - 2 = 0$ έχει συντελεστή διεύθυνσης $\lambda_{\zeta} = -\frac{1}{1} = -1$.
Άρα $\lambda_{KA}(-1) = -1 \iff \lambda_{KA} = 1$.
Επομένως η ευθεία $(KA)$ έχει εξίσωση:
\[
y - y_K = \lambda_{KA}(x - x_K) \iff y - 5 = 1(x - 3)
\]
\[
\iff y = x - 3 + 5 \iff y = x + 2 \iff x - y + 2 = 0.
\]
ii. Το σημείο $A$ είναι το σημείο τομής της ευθείας $\zeta$ με την ευθεία $KA$. Για τον προσδιορισμό των συντεταγμένων επιλύουμε το σύστημα των εξισώσεων των ευθειών $\zeta$ και $KA$.
\[
\left\{
\begin{array}{l}
x + y - 2 = 0 \\
x - y + 2 = 0
\end{array}
\right\}
\]
Προσθέτοντας τις δύο εξισώσεις κατά μέλη: $(x+y-2) + (x-y+2) = 0 \implies 2x = 0 \implies x = 0$.
Αντικαθιστώντας $x=0$ στην πρώτη εξίσωση: $0 + y - 2 = 0 \implies y = 2$.
Άρα το $A$ έχει συντεταγμένες $(0,2)$.

(γ) Η ευθεία $\varepsilon$ τέμνει τον κύκλο στα σημεία $M, \Lambda$ που είναι αντιδιαμετρικά. Επομένως $M\Lambda = 2\rho$.
Το ύψος του τριγώνου $A\Lambda M$ προς την $M\Lambda$ είναι ίσο με την απόσταση του σημείου $A$ από την ευθεία $\varepsilon$.
Είναι $\varepsilon: y = 2x - 1 \iff 2x - y - 1 = 0$, άρα $u = d(A, \varepsilon)$.
Το σημείο $A$ είναι $(0,2)$.
\[
u = \frac{|2 \cdot 0 - 2 - 1|}{\sqrt{2^2 + (-1)^2}} = \frac{|-3|}{\sqrt{4 + 1}} = \frac{3}{\sqrt{5}}.
\]
Οπότε το εμβαδόν του τριγώνου $A\Lambda M$ είναι:
\[
(A\Lambda M) = \frac{M\Lambda \cdot u}{2} = \frac{2\rho \cdot u}{2} = \rho \cdot u = 3\sqrt{2} \cdot \frac{3}{\sqrt{5}} = \frac{9\sqrt{2}}{\sqrt{5}} = \frac{9\sqrt{2}\sqrt{5}}{\sqrt{5}\sqrt{5}} = \frac{9\sqrt{10}}{5}.
\]
\end{lysh}